% Generated from project outputs. Do not edit by hand.
\begin{table}[!htbp]
\centering
\begin{threeparttable}
\caption{Labor income by subsample: 2023 moments}
\label{tab:annual-2023-y1-moments}
\small
\setlength{\tabcolsep}{4pt}
\renewcommand{\arraystretch}{1.15}
\begin{tabularx}{\linewidth}{@{}Xrrrrr@{}}
\toprule
Sample & N & Mean & SD & Min & Max \\
\midrule
General & 94 & 2,769.3 & 1,378.0 & 632.6 & 6,389.2 \\
Men & 94 & 3,120.0 & 1,473.6 & 732.7 & 7,857.7 \\
Women & 94 & 2,191.6 & 1,508.1 & 174.5 & 6,404.9 \\
Urban & 94 & 2,960.4 & 1,424.3 & 454.9 & 6,923.3 \\
Rural & 94 & 2,334.6 & 1,154.5 & 514.7 & 5,833.7 \\
Indigenous & 94 & 2,324.3 & 1,201.2 & 493.7 & 6,300.8 \\
Non-Indigenous & 94 & 2,964.5 & 1,403.3 & 487.4 & 6,804.0 \\
Bottom 99 percent & 94 & 2,630.7 & 1,202.0 & 632.6 & 5,531.0 \\
Bottom 90 percent & 94 & 2,174.4 & 795.5 & 603.6 & 4,455.5 \\
Bottom 50 percent & 94 & 1,180.3 & 379.8 & 254.0 & 2,296.8 \\
\bottomrule
\end{tabularx}
\begin{tablenotes}[flushleft]
\footnotesize
\item Monthly quetzales. Unweighted distribution of survey-weighted cohort means before transition matching. N counts cohort cells, not individuals. SD is dispersion, not a standard error. Subsamples overlap. Cumulative income definitions and treatment of missing components follow the merging scripts.
\end{tablenotes}
\end{threeparttable}
\end{table}
